\documentclass[11pt]{article}
\usepackage{amsmath,amssymb,amsthm}
\newtheorem{theorem}{Theorem}
\title{A short clock certificate can be expensive to obtain}
\begin{document}
\maketitle

One fixed admission question suffices to separate certificate size
from the work of selecting measurements. The certificate has length
\(m\), but the coordinator can need \(m^2\) queries to find it.
The gap does not require a later consumer that changes its mind.

\paragraph{Interface.}
There are \(m\) disjoint acquisition blocks, each with two episodes.
The first episode starts successful in the anchor, the second
failed. Each block has \(m\) separate clock sources relevant
only there. Use the pinned model in which their complete streams
have maximum residual \(9/4\), fixed private offsets in \([0,1]\),
and stamps shifted by one in the relevant block with flags
\(0,1\). They have two successful records in other blocks.
With the two anchor timestamps separated by two units and
the block bases separated by ten, the previous calculation
gives a feasible sorted join at every offset. A relevant
source can move its failure to the first episode precisely
when its offset lies in \([3/4,1]\).
An ideal query asks one named source for this historical bit.
All bit assignments are possible by choosing offsets zero
or one. Neither the source of the bit nor parallel read
latency is assumed: work counts total individual queries,
with exact decisions required in both outcomes.

The worst compatible failure count is \(m\) plus the number
of blocks with at least one affirmative bit. Use threshold
\(2m-1\), so the cap is violated iff every block has an
affirmative bit. Padding isolated, forced-safe episodes can
fix the cap fraction without changing the queries.

\begin{theorem}
For any \(m\ge1\), the shortest certificate on every input has
\(m\) answers; exact deterministic decision needs \(m^2\) in the
worst case. For zero-error randomized algorithms the optimal
worst-case expected number is \(m(m+1)/2\).
\end{theorem}
\begin{proof}
To certify violation, show one affirmative bit in each block.
Every certificate must show one in every block, or there
is a completion in which the omitted block is all negative.
To certify the cap, show a block with all answers negative.
Without such a block, every block can still have an affirmative
bit. Thus both minimal certificates cost \(m\).

The deterministic adversary answers negative on the first
\(m-1\) queries of each block and affirmative on the last.
Before all \(m^2\) queries, any unfinished block can be assigned
all negatives or a final affirmative, so both verdicts remain
possible. Reading all bits attains the bound.

For the randomized lower bound, consider positive inputs with
one affirmative per block, independently uniform among the
\(m\) locations. Every exact algorithm must find every affirmative.
Conditional on observed answers the remaining location in any
unresolved block is uniform, including when the queries are
interleaved with work on other blocks. Its expected rank in
the query order is \((m+1)/2\). Sum over blocks and average over
inputs to bound the worst expected work for any algorithm.

For the matching upper bound, uniformly permute the blocks and
the sources within each block. Stop reading a block at its
first affirmative and stop the procedure on finding an
all-negative block. If no block is all negative, each costs
at most \((m+1)/2\) in expectation. With \(z\ge1\) all-negative
blocks, the expected number of positive blocks before the first
negative is \((m-z)/(z+1)\), by symmetry of the \(z+1\) gaps
in the order. The expected total is at most
\[
     m+\frac{(m-z)(m+1)}{2(z+1)}
       \le \frac{m(m+1)}{2}.
\]
The last inequality follows by substituting
\(z=1\) for the largest
left-hand expression. This proves the value without an assumed
probability law on real clocks.
\end{proof}

\paragraph{What the test separates.}
At four blocks the certificate, exact deterministic worst work,
and zero-error worst expected work are respectively 4, 16 and 10.
A coordinator can simply refuse to admit an unresolved case;
the theorem is about deciding both branches, not about the least
cost of a safe refusal. If a source already supplies a trustworthy
aggregate per block, the coordinator reads fewer values, but the
proof must then account for evidence acquisition where the
aggregate was produced. Parallel queries may reduce wall time
without changing the work count. A recent polling response is
not itself the historical bit; the companion LinuxPTP and rate
model results explain which additional evidence would license it.

The replay independently evaluates the two-record timestamp
permutations for 530 Boolean fixtures. A decision-tree recurrence
checks deterministic depth through seven blocks, and direct
enumeration of randomized block orders checks the constructive
strategy through five. This is a scoped classical AND-of-OR
query model realized by the project's clock constraints, not a
claim of a new decision-tree lower-bound technique, measured
sensor cost, or full L1/R4 acceptance.
\end{document}
